\subsection{The annular pole and the two-wave estimate}
\label{sec:ex-annulus-two-wave}

L\"uscher and Weisz \cite[\S1]{LuscherWeiszCoupling} compute perturbative
coupling matching for the Wilson $\mathrm{SU}(N)$ action. The calculation
below identifies the logarithmic coefficient for the present real
Lie-algebra inner product and estimates the regulated finite-volume
response used in the scale iteration.


The selected one-step construction needs two facts about one finite-volume
response. Its logarithmic coefficient is an integral of the response over
four-dimensional momentum space, and that response is close to the continuum
quadratic pole on an annulus. This section gives the exact polar calculation
and expands the finite-volume two-channel estimate that supplies the two
radial bounds. The weighted-face estimates that feed it are proved in
Subsection~\ref{sec:ex-annulus-fleet-faces}; they are stated here at their
precise input interface.

\subsubsection{From the finite response to a momentum integral}

Let $E_4=(-\pi,\pi]^4$ be the fundamental momentum cube in Euclidean
coordinates, with normalized measure $c_{\mathrm{wt}}\,dx$ and
$c_{\mathrm{wt}}=(2\pi)^{-4}$. It has finite volume and contains the unit
ball. Two different objects occur, and the coefficient functional is
defined on the second.

A \emph{response kernel} $F$ has entries $F_{\xi\zeta}(p,x)$, indexed by
two marks (direction and color) and depending on an external momentum
$p$ and an internal momentum $x$. We assume that every entry is jointly
holomorphic, periodic and bounded on a common strip of positive width in
$(p,x)$. Its associated \emph{picture} is one half of the normalized
integral over the internal momentum,
\[
 \mathcal Y_{\xi\zeta}(p)=\frac12\,c_{\mathrm{wt}}\int_{E_4}
                         F_{\xi\zeta}(p,x)\,dx.
\]
For a two-leg picture $\mathcal Y$ and a color $a$, the coefficient
functional is twice the real part of the quadratic Taylor coefficient
along the axis $e_0$ of the component with both marks equal to
$(1,a)$:
\[
 \alpha_a(\mathcal Y)=2\operatorname{Re}\,[t^2]\,
                    \mathcal Y_{(1,a)(1,a)}(te_0)
 =\operatorname{Re}\frac{d^2}{dt^2}\mathcal Y_{(1,a)(1,a)}(te_0)
   \Big|_{t=0}.
\]
For the kernel put
\[
 \mathcal K_a(x)=\frac12\frac{d^2}{dt^2}F_{(1,a)(1,a)}(te_0,x)\Big|_{t=0},
 \qquad g_a(x)=\operatorname{Re}\mathcal K_a(x).
\]
Differentiation under the integral gives, for the picture $\mathcal Y$
associated with $F$,
\begin{equation}
 \alpha_a(\mathcal Y)=c_{\mathrm{wt}}\int_{E_4}g_a(x)\,dx.
 \label{eq:annulus-two-wave-estimates-display-AT1}
\end{equation}
The common analytic bound and Cauchy's estimate bound the second
$t$-derivative of the integrand uniformly in $x$, which justifies the
differentiation and shows that $g_a$ is bounded and continuous on the
closure of $E_4$. Equation \eqref{eq:annulus-two-wave-estimates-display-AT1} is an identity for the analytically
assembled response picture. The comparison of that picture with the
finite sampled pole test is a separate recording estimate with an
error term; it is not an equality and is not asserted by \eqref{eq:annulus-two-wave-estimates-display-AT1}.

For $x\in\mathbb R^4\setminus\{0\}$ let $r=|x|$ and define the gauge pole
\begin{equation}
 \Pi_G(x)=4C_G\left(
   \frac{x_1^2}{r^6}-4\frac{x_1^2x_0^2}{r^8}-\frac1{r^4}
 \right),
 \label{eq:annulus-two-wave-estimates-display-AT2}
\end{equation}
where the gauge group is simple and $C_G$ is its adjoint Casimir in the
normalization of the colour inner product on the Lie algebra. For a
group of unitary $r_G\times r_G$ matrices this inner product is
\[
 \langle u,v\rangle=-\frac1{r_G}\operatorname{Re}\operatorname{tr}(uv),
\]
the colour basis $(e_b)$ used below is orthonormal for it, and
\[
 -\operatorname{tr}(\operatorname{ad}u\,\operatorname{ad}v)
                =C_G\langle u,v\rangle,
 \qquad
 C_G=-\frac1{\dim\mathfrak g}\sum_b
       \operatorname{tr}\bigl((\operatorname{ad}e_b)^2\bigr).
\]
Simplicity is what makes the left side a multiple of the inner product.
The normalization matters for the constant $\lambda_G$ below: with the
trace form $-\operatorname{tr}(uv)$ in place of $\langle u,v\rangle$
the constant $C_G$ would be divided by $r_G$. For $SU(N)$ in its
defining representation, $C_G=2N^2$.
The spherical area of $S^3$ is $2\pi^2$. Rotational invariance gives
\begin{equation}
 \int_{S^3}\theta_1^2\,d\theta=\frac{2\pi^2}{4},\qquad
 \int_{S^3}\theta_1^2\theta_0^2\,d\theta=\frac{2\pi^2}{24}.
 \label{eq:annulus-two-wave-estimates-display-AT3}
\end{equation}
For the second identity, set $A=\int\theta_0^4$ and
$B=\int\theta_0^2\theta_1^2$. Coordinate permutations make these
independent of the distinct indices. From $\sum_i\theta_i^2=1$,
$4A+12B=2\pi^2$. Rotating the $0$-axis to
$(e_0+e_1)/\sqrt2$ gives $A=(2A+6B)/4$, hence $A=3B$ and
$B=2\pi^2/24$. The first identity follows by integrating
$\sum_i\theta_i^2=1$.

Consequently the angular mean of the parenthesis in \eqref{eq:annulus-two-wave-estimates-display-AT2} is
\[
 r^{-4}\left(\frac14-\frac4{24}-1\right)
       =-\frac{11}{12}r^{-4}.
\]
Let $0<\mu\leq1$. Since the Euclidean polar Jacobian is
$r^3dr\,d\theta$,
\[
 \int_{\mu<|x|\leq1}|x|^{-4}\,dx
       =2\pi^2\int_\mu^1\frac{dr}{r}
       =2\pi^2\log\mu^{-1}.
\]
Using $c_{\mathrm{wt}}=(2\pi)^{-4}$, we obtain the exact pole law
\begin{equation}
 c_{\mathrm{wt}}\int_{\mu<|x|\leq1}\Pi_G(x)\,dx
 =-\frac{11C_G}{24\pi^2}\log\mu^{-1}
 =-\lambda_G\log\mu^{-1},
 \qquad \lambda_G=\frac{11C_G}{24\pi^2}.
 \label{eq:annulus-two-wave-estimates-display-AT4}
\end{equation}
This fixes the sign and normalization of the logarithm; in particular the
pole is not defined only up to an angular constant.

\subsubsection{The finite-volume two-channel calculation}

We now state the two scale receipts used by the finite-volume construction.
Let $x\in[-\pi,\pi]^4$ be a centered momentum, $r=|x|$, and
$0<\mu\leq1$. The response kernel is the one of
Subsection~\ref{sec:ex-annulus-fleet-faces}: for a local record with color
$b$ and regulator $\mu$, it is one half of the two-leg figure of the
vector-channel total film minus that of the ghost-channel total film.
Write $\mathcal K(x)\in\mathbb C$ for its coefficient $\mathcal K_b(x)$
defined before \eqref{eq:annulus-two-wave-estimates-display-AT1}, and $g(x)=\operatorname{Re}\mathcal K(x)$. With
the faces $U,L,H$ of \eqref{eq:annulus-fleet-face-estimates-display-F5} in that section, formed in each of the two
channels with the same background $x$, the same mark $(1,b)$, the same
boost $e_0$ and the same regulator $\mu$, the six complex
second-derivative traces are
\[
\begin{array}{lll}
 D=U_{\rm vec}''(0),& d=L_{\rm vec}''(0),& d_0=H_{\rm vec}''(0),\\[2pt]
 Y=U_{\rm gh}''(0),& c=L_{\rm gh}''(0),& c_0=H_{\rm gh}''(0).
\end{array}
\]
The halo traces $c_0,d_0$ and the pole $\Pi_G$ are defined only for
$r>0$. The exact identities are
\begin{equation}\label{eq:ex-two-wave-identities}
 \mathcal K=\frac12\left(\frac12(D-d)-(Y-c)\right),\qquad
 \Pi_G=\frac12\left(c_0-\frac12d_0\right),
\end{equation}
\begin{equation}
 \mathcal K-\Pi_G=\frac12\left(\frac12D-Y\right)
       +\frac12\left((c-c_0)-\frac12(d-d_0)\right).
 \label{eq:annulus-two-wave-estimates-display-AT5b}
\end{equation}
The first identity holds at every centered $x$, including $x=0$: the
two-leg figure of a film is the trace of its upper face minus that of
its lower face, and the response kernel is one half of the vector
figure minus the ghost figure. The second is the free halo
identity, with marked direction and color fixed; it and \eqref{eq:annulus-two-wave-estimates-display-AT5b} hold for
$r>0$ and are used only on $\mu\leq r$. The third is their
algebraic difference; the real-valued pole $\Pi_G$ is embedded in $\mathbb C$
in these formulas. None of the three involves a torus side.

For clarity, the second identity follows directly from the two free
propagator numerators, that is, from the linear legs $w_b$ (ghost) and
$v_b$ (vector) of Subsection~\ref{sec:ex-annulus-fleet-faces} and the free
inverses $|k|^{-2}I$. In the selected unit color direction $e_b$, use the
normalization $\operatorname{tr}(\operatorname{ad}(e_b)^2)=-C_G$. Put
$\rho=|x|^2$, $c_{\rm col}=-C_G$, and $L=-4x_1^2$. Along the marked
displacement $t e_0$, the common denominator is
$\Delta(t)=\rho+2x_0t+t^2$. The ghost numerator is the scalar orbital
factor $L$ times the color trace, giving $Lc_{\rm col}/\rho$. For the
vector numerator, write $p=-2i x_1$ and let $A(t)$ have entries
$A_{01}(t)=-2it$, $A_{10}(t)=2it$, and all other entries zero. Then
$p^2=L$ and direct multiplication gives
\[
 \operatorname{tr}\bigl[(pI+A(-t))(pI+A(t))\bigr]
       =4p^2-8t^2=4L-8t^2.
\]
Multiplying by the same color trace and scalar denominator factor gives
the vector numerator $\rho^{-1}(4L-8t^2)c_{\rm col}$. Differentiating
the quotients twice at $t=0$ gives
\[
 \left(\frac{Lc_{\rm col}/\rho}{\Delta}\right)''_{t=0}
       =-\frac{2Lc_{\rm col}}{\rho^3}
           +\frac{8Lc_{\rm col}x_0^2}{\rho^4},\qquad
 \left(\frac{\rho^{-1}(4L-8t^2)c_{\rm col}}{\Delta}\right)''_{t=0}
       =-\frac{16c_{\rm col}}{\rho^2}
         -\frac{8Lc_{\rm col}}{\rho^3}
         +\frac{32Lc_{\rm col}x_0^2}{\rho^4}.
\]
Thus the orbital trace minus one half of the vector trace is
\[
 \frac{8c_{\rm col}}{\rho^2}
       +\frac{2Lc_{\rm col}}{\rho^3}
       -\frac{8Lc_{\rm col}x_0^2}{\rho^4}
 =8C_G\left(\frac{x_1^2}{r^6}
       -4\frac{x_1^2x_0^2}{r^8}-\frac1{r^4}\right)
 =2\Pi_G(x),
\]
which proves the halo identity in \eqref{eq:ex-two-wave-identities}. \eqref{eq:annulus-two-wave-estimates-display-AT5b} is then obtained by
subtracting the just-computed free identity from the first identity and
collecting terms.

The first family of estimates is the all-centered-record receipt, with
$T_1=r+\mu$ and $R_0=1+2\pi$:
\begin{equation}
 0\leq r,\quad T_1=r+\mu\leq R_0,\qquad
 |\mathcal K|\leq A_1T_1^{-4}.
 \label{eq:annulus-two-wave-estimates-display-AT6}
\end{equation}
The second is the shell receipt, available only on $\mu\leq r\leq1$ and
using $T_2=r$:
\begin{equation}
 |D|\leq u r^{-2}+v\mu^2r^{-6},\qquad
 |Y|\leq u_{\mathrm{ref}}r^{-2}+v_{\mathrm{ref}}\mu^2r^{-6},
 \label{eq:annulus-two-wave-estimates-display-AT7}
\end{equation}
\begin{equation}
 |c-c_0|\leq A_0(r+\mu^2r^{-2})r^{-4},\qquad
 |d-d_0|\leq B_0(r+\mu^2r^{-2})r^{-4}.
 \label{eq:annulus-two-wave-estimates-display-AT8}
\end{equation}
All constants $A_1,u,v,u_{\mathrm{ref}},v_{\mathrm{ref}},A_0,B_0$ are nonnegative and
are chosen before the local record: they do not depend on its regulator,
color, bed, clip or strip width. The subscript ${\rm ref}$ marks the
ghost channel. These two receipts are kept
separate: the all-centered estimate uses scale $r+\mu$, while the shell
comparison uses scale $r$. They are the conclusions \eqref{eq:annulus-fleet-face-estimates-display-F7} and \eqref{eq:annulus-fleet-face-estimates-display-F8} of
Proposition~\ref{prop:ex-fleet-face-bounds}, applied in the two channels
to the family of all centered records with the prices derived there
(Lemma~\ref{lem:ex-row-to-floe} and the trajectory-price paragraph):
\eqref{eq:annulus-two-wave-estimates-display-AT6} follows from \eqref{eq:annulus-fleet-face-estimates-display-F7} at scale $r+\mu$ and the first identity in
\eqref{eq:ex-two-wave-identities}, \eqref{eq:annulus-two-wave-estimates-display-AT7} is the upper bound in \eqref{eq:annulus-fleet-face-estimates-display-F7} at scale $r$, and \eqref{eq:annulus-two-wave-estimates-display-AT8} is \eqref{eq:annulus-fleet-face-estimates-display-F8}
at scale $r$.

\begin{proposition}[Two wave bounds from the finite response]
\label{prop:ex-two-wave}
Suppose \eqref{eq:ex-two-wave-identities}--\eqref{eq:annulus-two-wave-estimates-display-AT8} hold uniformly. Then
\begin{equation}
 |g(x)|\leq \frac{A_1}{(r+\mu)^4},\qquad
 |g(x)-\Pi_G(x)|\leq \frac{A_2}{r^3}
                     +\frac{A_2\mu^2}{r^6}
 \quad(\mu\leq r\leq1),
 \label{eq:annulus-two-wave-estimates-display-AT9}
\end{equation}
where
\[
 A_2=\max\{u+u_{\mathrm{ref}}+A_0+B_0,
                    v+v_{\mathrm{ref}}+A_0+B_0\}.
\]
\end{proposition}
\begin{proof}
The first estimate follows from \eqref{eq:annulus-two-wave-estimates-display-AT6} and $|\operatorname{Re}z|\leq|z|$.
On the annulus, use \eqref{eq:annulus-two-wave-estimates-display-AT5b} and \eqref{eq:annulus-two-wave-estimates-display-AT7}--\eqref{eq:annulus-two-wave-estimates-display-AT8}:
\begin{equation}
 |\mathcal K-\Pi_G|
 \leq (u+u_{\mathrm{ref}})r^{-2}
      +(v+v_{\mathrm{ref}})\mu^2r^{-6}
 +(A_0+B_0)(r^{-3}+\mu^2r^{-6}).
 \label{eq:annulus-two-wave-estimates-display-AT10}
\end{equation}
Here $\tfrac12$ factors are bounded by one, and
$(r+\mu^2r^{-2})r^{-4}=r^{-3}+\mu^2r^{-6}$. Since $r\leq1$,
$r^{-2}\leq r^{-3}$. The displayed $A_2$ bounds \eqref{eq:annulus-two-wave-estimates-display-AT10} by
$A_2(r^{-3}+\mu^2r^{-6})$. Taking real parts gives
$|g-\Pi_G|\leq|\mathcal K-\Pi_G|$, proving \eqref{eq:annulus-two-wave-estimates-display-AT9}.
\end{proof}

\subsubsection{The two applications of the finite-volume receipts}

The common family of local records is used twice, with different scale
parameters. Each record supplies a centered momentum $x$, a positive
regulator $\mu\leq1$, and the same complex response coefficient
$\mathcal K(x)$. The centered-box condition $|x_j|\leq\pi$ gives
$|x|\leq2\pi$.

For the all-centered estimate, use scale $T_1=|x|+\mu$. The whole-frame
bound from the two channels is
\[
 |D-d|\leq A_{\rm int}T_1^{-4},\qquad
 |Y-c|\leq A_{\rm ref}T_1^{-4}.
\]
The exact first identity in \eqref{eq:ex-two-wave-identities} then gives
$|\mathcal K|\leq(A_{\rm int}/4+A_{\rm ref}/2)T_1^{-4}$. This is the
first radial estimate; it applies also at $x=0$.

For the annular comparison, restrict the same records to
$\mu\leq|x|$ and use the second scale $T_2=|x|$. The upper face of the
vector channel is split
into a principal term and a tadpole term; their charged bounds are
$uT_2^{-2}$ and $v\mu^2T_2^{-6}$. The same split for the ghost channel
has bounds $u_{\mathrm{ref}}T_2^{-2}$ and
$v_{\mathrm{ref}}\mu^2T_2^{-6}$. The lower-channel
remainders after subtracting their respective halos are bounded by
$A_0(T_2+\mu^2/T_2^2)T_2^{-4}$ and
$B_0(T_2+\mu^2/T_2^2)T_2^{-4}$. The free quotient calculation gives
$\Pi_G=\frac12(c_0-\frac12d_0)$, so these four estimates feed exactly
into \eqref{eq:annulus-two-wave-estimates-display-AT5b} and yield \eqref{eq:annulus-two-wave-estimates-display-AT7}--\eqref{eq:annulus-two-wave-estimates-display-AT10}.

These scale receipts are thus two separate applications of the
finite-volume receipts to the common family of local records: the first
applies the frame estimate to every centered record, while the second
restricts to the subfamily $\mu\leq|x|$ before applying the shell
estimate. They use the same response
coefficient and common constants. The annulus proof therefore does not
interchange a limit with a volume-dependent integral. The comparison of
this analytic coefficient with the finite sampled pole test is a
separate recording estimate, with an error equal to the site count
times a quantity tending to zero; exact equality with the finite
sampled coefficient is not asserted, and nothing of the kind is
inferred from \eqref{eq:annulus-two-wave-estimates-display-AT1}.

\subparagraph{Scope.}
This section proves the normalized pole integral, the halo identity and
the algebraic conversion from the two scale receipts to the radial
bounds. The receipts themselves are proved in
Subsection~\ref{sec:ex-annulus-fleet-faces} for the family of all centered
records. The recording estimate is proved in
Subsection~\ref{sec:ex-annulus-recording}. The order of choices there
differs from the one for the radial constants: the fee strip of the
response and the constants $A_1,A_2$ are fixed before the blocking side
$L$, whereas the strip width $z_L$ and the bound $K_L$ of the regulated
recording kernel are obtained after $L$ is fixed, with $\mu=L^{-1}$,
and may depend on $L$. For fixed $L$ the recording error is the site
count times $\mathfrak r(z_L,K_L)e^{-z_L\ell/4}$, in the notation of
that section, which tends to zero along every sequence of sides
$\ell\to\infty$.
